% =====================================================================
%  Proof-Decomposition Granularity for Small Theorem Provers
%  LIVING DRAFT — incomplete by design. Fill incrementally.
%  Target venue: ICLR 2027.  Swap this preamble for the official
%  iclr20XX_conference.sty when the style file is released; the
%  \documentclass/article setup below is a placeholder so the draft
%  compiles standalone in the meantime.
% =====================================================================
\documentclass[11pt]{article}
\usepackage[margin=1in]{geometry}
\usepackage{amsmath,amssymb,amsthm,mathtools}
\usepackage{hyperref}
\usepackage{booktabs}
\usepackage{xcolor}
\usepackage{natbib}

\newcommand{\TODO}[1]{\textcolor{red}{\textbf{[TODO: #1]}}}
\newcommand{\PLACEHOLDER}[1]{\textcolor{blue}{\emph{[placeholder: #1]}}}
\newtheorem{definition}{Definition}
\newtheorem{proposition}{Proposition}

\title{Proof-Decomposition Granularity is a Controllable\\
       Training-Data Axis for Small Theorem Provers}
\author{\TODO{authors} }
\date{Draft compiled \today}

\begin{document}
\maketitle

\begin{abstract}
Long formal proofs are hard for small models for the same reason long-horizon
control is hard for learned world models: compounding step-wise error and a
search space that grows exponentially in the horizon. In control, this is
mitigated by \emph{hierarchical planning} across temporal scales
\citep{zhang2026hierarchical}. We ask the analogous question for theorem
proving: \emph{does the granularity at which a proof is decomposed into
sub-results change what a small prover can learn?} We formalize a proof as a
node in the acyclic dependency graph of a library, made finite by a bounded
\emph{expansion frontier}, and define a \emph{decomposition} as a
\emph{section} --- a maximal antichain (cut) of the resulting proof poset. A
section induces a Lean-verified rewrite of the proof; semantics preservation is
not assumed but \emph{checked}. We compare a family of systematic
section-selection policies as a controlled, \emph{within-theorem}
training-data axis, fine-tune small ($\leq$1.5\,B) provers under a fixed token
budget, and measure held-out \texttt{pass@}$k$ on miniF2F and ProofNet.
\PLACEHOLDER{headline result: an intermediate section granularity dominates the
flat and maximally-decomposed extremes, with the effect growing as model size
decreases.}
\end{abstract}

\section{Introduction}
\label{sec:intro}

\paragraph{The horizon problem, twice.} \citet{zhang2026hierarchical} show that
model-predictive control with a single-scale learned world model fails on
long-horizon tasks because prediction errors compound and the planning search
space grows exponentially; planning hierarchically across temporal scales
recovers long-horizon performance, and the advantage of the hierarchical
planner over the flat one \emph{grows with the horizon}. Autoregressive
theorem provers face a structurally identical problem: tactic-level error
compounds over proof length and proof search is exponential in depth. This
motivates a proof-theoretic analog of temporal abstraction --- decomposing a
proof through intermediate results --- and the central empirical question of
this paper: \emph{which decomposition?}

\paragraph{Contributions.}
\begin{enumerate}
  \item A formalization of proof decomposition as a \emph{section} (maximal
        antichain) of the bounded-unfolded library dependency DAG
        (\S\ref{sec:formalism}), which strictly generalizes
        \texttt{have}-tree decompositions.
  \item A \emph{generate-then-verify} pipeline: every decomposition is
        reconstructed into Lean and machine-checked against the original
        theorem with an axiom audit, so semantics preservation is a verified
        property, not an assumption (\S\ref{sec:method}).
  \item A controlled, within-theorem study of section-selection policies as a
        training-data axis for small provers, with explicit decorrelation of
        granularity from sequence length and difficulty
        (\S\ref{sec:protocol}).
  \item \PLACEHOLDER{empirical finding + the granularity $\times$ model-size
        interaction.}
\end{enumerate}

\section{Formalism: proofs, frontiers, sections}
\label{sec:formalism}

\paragraph{The proof DAG.} Let $V$ be the named declarations of the library.
For $d\in V$, its elaborated proof term references constants $C(d)\subseteq V$;
edges $d\to e$ for $e\in C(d)$ form an acyclic graph $G=(V,\to)$ (proofs are
well-founded). Within one declaration's tactic block, \emph{internal nodes}
(\texttt{have}, \texttt{suffices}, \texttt{calc} steps, \dots) carry a
typed statement and a justification, forming a tree hanging off $d$.

\paragraph{Bounded unfolding.} Unfolding $G$ from a root $t$ into a tree is
super-exponential and bottoms out only at primitive recursors. We therefore
define the proof tree $T(t,F)$ relative to an \emph{expansion frontier}
$F\subseteq V$: recurse into a child $e$ iff $e\notin F$; nodes in $F$ and
primitives are opaque leaves. $F=V$ recovers the in-proof \texttt{have}-tree
as a special case.

\begin{definition}[Section]
\label{def:section}
A \emph{section} of $T(t,F)$ with order $\preceq$ is a set $S$ that is
(i)~an antichain (no two elements $\preceq$-comparable) and (ii)~maximal:
every node of $T$ is $\preceq$-comparable to some $s\in S$. Equivalently,
every root-to-leaf path meets $S$ exactly once.
\end{definition}

A section partitions $T$ into the \emph{upper part} (the proof of $t$ to be
produced, using $S$ as named lemmas) and the subtrees below each $s\in S$.
The extremes are $S=\{t\}$ (flat, no decomposition) and $S=\mathrm{leaves}(T)$
(maximal decomposition at $F$).

\section{Method: policies and verified reconstruction}
\label{sec:method}

\paragraph{Section-selection policies.} The number of maximal antichains is
exponential; we compare a fixed family of policies $\pi:(t,F)\mapsto S$:
$\pi_{\mathrm{root}}$, $\pi_{\mathrm{leaf}}$ (controls);
$\pi_{\mathrm{depth}(k)}$, $\pi_{\mathrm{size}(\tau)}$,
$\pi_{\mathrm{premise}(p)}$, $\pi_{\mathrm{balanced}}$, $\pi_{\mathrm{named}}$
(systematic interior); and a learned predictor $\pi_{\mathrm{learned}}$ as a
hierarchical-planning analog of \citet{zhang2026hierarchical} (Variant~B).

\paragraph{Verified reconstruction.} Each $(t,F,S)$ is rendered to a single
Lean artifact (section nodes as \texttt{have}/auxiliary lemmas) and checked:
it must compile against $t$'s original statement and \texttt{\#print axioms}
must not exceed the original's axiom set (no \texttt{sorry}, no smuggled
axioms). Non-verifying decompositions are discarded; reconstruction yield is
reported. All verification uses a single frozen toolchain --- Lean
\texttt{v4.26.0}, Mathlib \texttt{2df2f015\ldots}, \texttt{repl}
\texttt{a9a6f5ba\ldots} --- shared byte-identically by extraction (a fork of
ntp-toolkit \citep{hu2024minictx} at its verified v4.26.0 state), rewriting,
and evaluation, eliminating toolchain skew between the corpus and the
verifier.

\section{Experimental protocol}
\label{sec:protocol}

\paragraph{Within-theorem design \& decorrelation.} The set of root theorems
is identical across all policy corpora; only articulation differs. Granularity
is decorrelated from sequence length and difficulty (matched token budgets,
length as a regression covariate, difficulty strata, and the
length-controlled policy $\pi_{\mathrm{size}}$).

\paragraph{Models, training, evaluation.} Primary base
Kimina-Prover-Distill-1.5B \citep{kimina2025}; controls Qwen2.5-Math-1.5B
and Qwen2.5-0.5B;
QLoRA on a single T4 at fixed token budget. Held-out \texttt{pass@}1/8 on
miniF2F \citep{zheng2022minif2f} (and the corrected miniF2F\_v2
\citep{minif2f_v2}) and ProofNet \citep{azerbayev2023proofnet}, all checked by
the same Lean harness. Full configuration and decision gates: see the
companion protocol \texttt{EXPERIMENT.md}.

\section{Related work}
\label{sec:related}

\paragraph{Hierarchical planning and abstraction.} Our framing is borrowed
directly from hierarchical planning with multi-scale world models
\citep{zhang2026hierarchical}: their temporal abstraction is, in the proof
setting, a section of the dependency DAG, and their finding that intermediate
abstraction beats single-level planning on long horizons is the control-domain
analog of our central hypothesis. We test that principle where a verifier
provides ground truth.

\paragraph{Sketch- and subgoal-based proving.} Draft--Sketch--Prove
\citep{jiang2023draft}, LEGO-Prover \citep{wang2024lego}, Lean-STaR
\citep{lin2024leanstar}, and subgoal decomposition in DeepSeek-Prover
\citep{xin2024deepseekprover,ren2025deepseekproverv2} \emph{generate}
sketches or subgoals at inference. ProofAug \citep{liu2025proofaug} analyzes proof structure and
inserts automation at multiple granularities, also at inference time, and
likewise improves a 1.5\,B prover. \textbf{Distinction:} we vary decomposition
granularity as a \emph{controlled training-data axis} over the \emph{same}
verified proofs, not as an inference-time procedure.

\paragraph{Extraction and benchmarks.} We build on the miniCTX/ntp-toolkit
extraction stack \citep{hu2024minictx} and LeanDojo \citep{yang2023leandojo},
and report on miniF2F \citep{zheng2022minif2f,minif2f_v2} and ProofNet
\citep{azerbayev2023proofnet}.

\section{Results}
\label{sec:results}
\PLACEHOLDER{Stage 0.5 pilot, then Stage 3 main table: held-out pass@k vs.
section-policy granularity, per model size; verified-reconstruction yield;
length-controlled regression; per-Mathlib-area ablation.}

\section{Conclusion}
\PLACEHOLDER{}

\bibliographystyle{plainnat}
\bibliography{refs}

\end{document}
